% ACL-style draft for the NarrativeSimilarity project.
% Compile from this directory with: latexmk -pdf main.tex

\documentclass[11pt]{article}

% For camera-ready, remove the [review] option.
\usepackage[review]{acl}

\usepackage{times}
\usepackage{latexsym}
\usepackage[T1]{fontenc}
\usepackage[utf8]{inputenc}
\usepackage{microtype}
\usepackage{inconsolata}
\usepackage{graphicx}
\usepackage{booktabs}
\usepackage{amsmath}
\usepackage{amssymb}
\usepackage{multirow}
\usepackage{xcolor}
\usepackage{url}

\title{Learning Structure-Aware Story Embeddings for Narrative Similarity}

\author{First Author \\
  Affiliation \\
  \texttt{email@example.com} \\
  \And
  Second Author \\
  Affiliation \\
  \texttt{email@example.com}}

\begin{document}
\maketitle

\begin{abstract}
We study the problem of learning story embeddings that capture structural similarity: similarity in what happens, how events unfold, and how narrative roles progress, rather than only topical or lexical overlap. We combine a human-calibrated structural similarity model with pseudo-labeled story pairs to fine-tune embedding models and evaluate them across pairwise, retrieval, ranking, synthetic, and external benchmark tasks.
\end{abstract}

\section{Introduction}

Narrative similarity is often treated as a semantic matching problem, where models prioritize lexical overlap, topical closeness, or stylistic resemblance. However, many downstream tasks require sensitivity to \emph{what happens} in a story and \emph{how events unfold} over time. Two narratives may use different wording and details yet still share the same underlying event progression, while semantically similar texts may diverge in causal and temporal structure.

This distinction matters for practical systems. In retrieval-augmented generation (RAG), a retriever that captures narrative structure can return context that is behaviorally relevant rather than merely topically related. In recommendation systems, users often care about plot dynamics (e.g., setup, escalation, reversal, resolution), not only keywords or genres. As large language models continue to improve, direct LLM-based scoring is increasingly effective but still expensive to scale across large corpora. A strong embedding model that approximates structure-aware judgments offers a lower-latency, lower-cost alternative for large-scale indexing, retrieval, and ranking.

In this work, we focus on learning story embeddings that are explicitly sensitive to \emph{structural similarity}. Our goal is not simply to improve generic semantic similarity, which is already well served by strong off-the-shelf embedding models. Instead, we ask whether an embedding model can be trained to place stories close together when they share event progression, temporal organization, and narrative dynamics, even when they differ in topic, entities, wording, or style.

A central challenge is that structural similarity is not a standard supervised target. To ground the task, we run a user study that asks how people perceive structural similarity between narratives. This study is important because it determines what our training signal should emphasize: not just whether two stories discuss related content, but whether their event sequences align in a way that humans recognize as structurally similar. We use these judgments to calibrate a structural similarity model that combines event alignment and aligned-event semantics.

The resulting calibrated model acts as a teacher for embedding training. We apply it to larger ASQ-derived story-pair data to obtain pseudo-labels, then fine-tune neural embedding models so cosine similarity reflects the teacher's structure-aware scores. We evaluate the resulting embeddings across several complementary datasets and task formats, including pair separation, retrieval, ranking, synthetic anchor comparisons, and external narrative-similarity benchmarks. This lets us test whether structural supervision transfers beyond the pairwise training setup into retrieval and ranking settings that better resemble practical applications.

\section{Related Work}

Prior work on semantic textual similarity, narrative understanding, event extraction, story schemas, and dense retrieval provides the foundation for this project. Semantic embedding models are effective for topical and meaning-based matching, but narrative structure requires attention to event order, role alignment, and plot progression.

This work is also related to learning from weak supervision and pseudo-labels, where a smaller calibrated model or annotation pipeline is used to generate training signals for a larger dataset. We position structural story embedding as a bridge between explicit narrative modeling and scalable neural retrieval.

\section{Datasets}

\subsection{Training Data: ASQ Pair Set}

We use ASQ-derived story pairs as the primary unlabeled pool for pseudo-label generation and embedding training. Pair-level structural scores are produced by the structural teacher model after obtaining alignments, yielding scalable supervision without requiring dense human annotations for all pairs.

\subsection{Evaluation Data}

\subsubsection{Movie Remakes}

Movie Remakes provides naturally paired narratives that describe similar underlying plots with variation in style and detail. We construct balanced evaluation sets with positive pairs (retellings/remakes) and random negatives, and evaluate both pairwise similarity separation and retrieval behavior.

\subsubsection{Tell Me Again}

Tell Me Again includes multiple retellings of the same source stories. For evaluation stability, we restrict sampling to English stories and build both pairwise and retrieval-style test sets. This dataset tests whether the model captures shared event progression despite paraphrastic variation.

\subsubsection{Synthetic Stories (v1/v2)}

We also evaluate on controlled synthetic narratives. Each theme includes core events and three generated stories ($\text{story}_1,\text{story}_2,\text{story}_3$). Data generation explicitly enforces chronological coherence and event-level separability. In \texttt{synthetic\_stories\_v2}, generation instructions specify that:
\begin{itemize}
    \item all stories cover the same scenario and core events,
    \item one event per sentence is preferred,
    \item direct phrase copying is discouraged,
    \item \texttt{story\_2} is intended to be most similar to \texttt{story\_1}, while \texttt{story\_3} introduces stronger structural differences.
\end{itemize}

This setup provides a controlled benchmark for anchor-based structural ranking.

\subsubsection{SemEval 2026 Task 4}

We include the development split of SemEval 2026 Task 4 as an external triplet-style evaluation. Each example contains an anchor text and two candidates, \texttt{text\_a} and \texttt{text\_b}, with a binary label indicating whether \texttt{text\_a} is closer to the anchor. This format directly tests whether cosine similarity in the embedding space selects the candidate that is judged closer to the anchor.

\subsubsection{SemEval 2022 Task 8}

We also evaluate on SemEval 2022 Task 8, a multilingual news article similarity benchmark annotated along several dimensions, including narrative similarity~\cite{chen2022semeval}. Because the released data primarily contains URLs and annotation scores, we construct an English--English subset with scraped article text and retain the \textsc{NAR} score as the human narrative-similarity target. This evaluation tests whether embedding cosine similarity correlates with human judgments of narrative schema similarity in news articles.

\section{Method}

Modeling structural similarity between narratives is challenging because there is no formal or universally accepted definition of how humans perceive two stories as being structurally similar. While prior work on narrative similarity often focuses on topical overlap or full-text semantic similarity, structural similarity depends more specifically on how events unfold, align, and relate across stories. To better capture this notion, we first conduct a human study on pairs of personal narratives sampled from the RTN dataset. Participants rate the structural similarity of story pairs, allowing us to obtain direct human judgments about perceived narrative structure.

Using these annotations, we train a structure-aware scoring model that combines two complementary components: (1) alignment similarity, which measures consistency in event ordering and event coverage across stories, and (2) semantic similarity between aligned events. The resulting model serves as a calibrated approximation of human structural similarity judgments. In the remainder of the paper, we treat this human-aligned scorer as a teacher model capable of assigning structural similarity scores to unseen story pairs.

Next, we scale beyond the relatively small human-annotated dataset by applying the trained structural scorer to a large collection of approximately 10,000 randomly constructed story pairs from the ASQ dataset. This process generates pseudo-labels representing estimated structural similarity scores between stories. These pseudo-labeled pairs are then used to train a neural embedding model that learns to encode stories according to structural similarity rather than surface-level semantic overlap or topic similarity. As a result, stories with similar narrative progression and event structure are mapped closer together in embedding space, even when their content differs substantially.

\subsection{Human Study on Structural Similarity}

We first conduct a human annotation study to obtain grounded similarity judgments. Using personal narratives from the RTN corpus~\cite{saldias2020exploring}, we construct pairs of stories and ask participants to rate their similarity under two conditions: (1) based on the full stories, and (2) based only on extracted event sequences. This within-subjects design allows us to compare holistic narrative similarity against structure-focused similarity judgments.

To represent narrative structure, we extract the main events from each story using GPT-5 (temperature = 0). We define an event as a singular, temporally grounded occurrence in the narrative, excluding habitual actions, background descriptions, and general states. The model is instructed to preserve chronological order and return a structured list of events for each story. We retain up to six events per narrative, which serve as the structural representation used throughout the study.

To ensure construction of informative evaluation pairs, we compute two independent similarity measures across all stories: (1) full-text semantic similarity using embeddings of the complete narratives, and (2) event-based similarity using embeddings of the extracted event sequences. Embeddings are generated with the \texttt{all-mpnet-base-v2} SentenceTransformer model~\cite{reimers2019sentence}. We then divide both similarity distributions into high and low regions using percentile-based thresholds and construct four pairing conditions: High--High, Low--Low, High--Low, and Low--High, corresponding to combinations of high or low full-text similarity and event similarity. This design enables us to isolate cases where structural similarity and surface-level semantic similarity either align or diverge.

Finally, we sample 25 story pairs from each condition, yielding a balanced set of 100 narrative pairs for annotation. The resulting human judgments form the supervision signal used to train and calibrate our structural similarity model, which later serves as the teacher model for large-scale pseudo-label generation and embedding training.

\subsection{Structural Similarity Model}

Our structural scorer combines two components over aligned event sequences:
\begin{enumerate}
    \item \textbf{Alignment similarity}: captures ordering consistency and coverage of aligned events.
    \item \textbf{Semantic similarity}: captures meaning similarity within aligned event pairs.
\end{enumerate}

Given two event sequences $A$ and $B$, and an alignment set $\mathcal{M}$ (possibly many-to-many), we define:
\begin{itemize}
    \item $\text{cost}_{\text{skip}}$: number of unaligned events in both stories.
    \item $\text{cost}_{\text{reorder}}$: number of inversions induced by aligned event order.
\end{itemize}

Let $e_A=|A|$, $e_B=|B|$, and $n=|\mathcal{M}_{\text{pairs}}|$ be the number of expanded aligned pairs. The normalized alignment distance is:
\begin{equation}
D_{\text{align}} = \frac{\text{cost}_{\text{reorder}} + \text{cost}_{\text{skip}}}{\frac{n(n-1)}{2} + (e_A + e_B)}.
\end{equation}

Alignment similarity is:
\begin{equation}
S_{\text{align}} = 1 - D_{\text{align}}.
\end{equation}

For semantic distance, we embed aligned event texts and compute pairwise cosine distance in $[0,1]$:
\begin{equation}
D_{\text{sem}} = \frac{1}{n}\sum_{(i,j)\in\mathcal{M}_{\text{pairs}}} \frac{1-\cos(\mathbf{u}_i,\mathbf{v}_j)}{2},
\end{equation}
with semantic similarity:
\begin{equation}
S_{\text{sem}} = 1 - D_{\text{sem}}.
\end{equation}

The final structural score is a calibrated linear model fit on human survey ratings:
\begin{equation}
\hat{y} = c + \alpha S_{\text{align}} + \beta S_{\text{sem}}.
\end{equation}

The learned parameters used in our implementation are:
\begin{equation}
(c,\alpha,\beta) = (1.5830,\ 0.2021,\ 0.9339).
\end{equation}

\subsection{Pseudo-Labeling and Embedding Training}

We use the structural scorer above as a teacher to generate pseudo-labels on a larger unlabeled pool of story pairs (ASQ-based pair construction). These pseudo-labels supervise embedding training so that the student model learns to approximate structure-aware similarity directly from full stories.

Our embedding backbone is \texttt{intfloat/e5-mistral-7b-instruct}, fine-tuned with LoRA. We train with a contrastive MSE objective that regresses pairwise similarity behavior toward the teacher signal (using the pseudo-labeled structure score framework). Checkpoints are saved periodically (every 2 epochs), enabling trajectory analysis across training.

\section{Evaluation}

We evaluate whether embedding cosine similarity captures structural story similarity across several tasks.

\textbf{Pairwise similarity.}
Each example contains two stories and a binary label indicating whether the pair is positive or negative. We report the mean cosine similarity for positive and negative pairs, and the gap between them.

\textbf{Retrieval.}
Each query story is paired with five candidate stories, one positive and four negatives. The model ranks candidates by cosine similarity to the query. Accuracy is the fraction of examples for which the positive candidate receives the highest similarity.

\textbf{Ranking.}
Each query story has five candidates with a full relevance ranking. We evaluate predicted rankings using normalized discounted cumulative gain (nDCG), which rewards placing more relevant candidates near the top.

\textbf{Synthetic anchor evaluation.}
For synthetic stories, one story is treated as an anchor and the model must identify which of two alternatives is more structurally similar according to the structural similarity score.

\textbf{SemEval evaluation.}
For SemEval 2026-style triplets, the model selects whether text A or text B is closer to the anchor. For SemEval 2022 Task 8, we compute rank correlation between embedding cosine similarity and the human narrative similarity score.

\section{Experiments}

We fine-tune embedding models on pseudo-labeled ASQ story pairs. The main model family includes E5/Mistral-style instruction embeddings and BGE-style encoder embeddings. Experiments compare checkpoints across epochs and evaluate whether training improves pairwise separation, retrieval accuracy, ranking quality, synthetic preference accuracy, and correlation with human narrative similarity scores. We also compare GPT-5 structural ratings against the calibrated structural scorer to assess how closely an LLM-as-judge setup matches the teacher signal used for pseudo-labeling.

\subsection{Training Procedure}

Our primary embedding model fine-tunes \texttt{intfloat/e5-mistral-7b-instruct} on ASQ-derived story pairs labeled by the structural teacher model. Each training example consists of two full story texts and a teacher-produced structural similarity score. We normalize teacher scores to $[0,1]$ and train the embedding model so that cosine similarity between the two story embeddings matches the normalized teacher score:
\begin{equation}
\mathcal{L}_{\text{MSE}} =
\left(\cos(f_\theta(x_i), f_\theta(x_j)) - \tilde{s}_{ij}\right)^2,
\end{equation}
where $f_\theta$ is the embedding model and $\tilde{s}_{ij}$ is the normalized structural pseudo-label for pair $(i,j)$. This contrastive regression objective preserves the pairwise nature of the teacher signal while producing reusable embeddings for retrieval and ranking.

For the E5-Mistral runs, stories are formatted with the E5 passage prefix, truncated to a maximum length of 1024 tokens, and trained with bfloat16 precision when available. We use AdamW with learning rate $1\times10^{-4}$, weight decay 0.01, linear warmup ratio 0.03, maximum gradient norm 1.0, batch size 8, and gradient accumulation over 8 steps. The training split uses 90\% of pseudo-labeled pairs for training and 10\% for validation with seed 42. We train for 10 epochs and save epoch checkpoints for downstream evaluation.

To make fine-tuning feasible for the large E5-Mistral backbone, we use LoRA adapters with rank 16, $\alpha=32$, and dropout 0.05 over the attention and feed-forward projection modules. We also implement a BGE-base training variant. For BGE, we use \texttt{BAAI/bge-base-en-v1.5} with maximum length 512 and no LoRA adapters; the current BGE configuration supports InfoNCE training with one anchor, one positive, and two negatives sampled from teacher-score bins. Unless otherwise stated, the main reported checkpoint comparison focuses on the E5-Mistral contrastive MSE setting, because it directly regresses to the calibrated structural score.

\subsection{Baselines}

We compare fine-tuned models against embedding baselines that represent general-purpose semantic encoders, retrieval-specialized encoders, and narrative-specific story representations. We also include a direct GPT-5 judge baseline where prompt-based evaluation is available.

\paragraph{RoBERTa-base.}
RoBERTa-base is a strong Transformer encoder trained with a robustly optimized masked-language-modeling objective~\cite{liu2019roberta}. It is not trained specifically for retrieval or narrative similarity, so it provides a useful general-purpose language-model baseline for testing whether structural evaluation requires more than contextual token representations.

\paragraph{BGE-base.}
We use \texttt{BAAI/bge-base-en-v1.5} as a retrieval-oriented embedding baseline. BGE models are designed for dense text embedding and retrieval, and the v1.5 release is intended to improve similarity-score behavior while preserving strong retrieval performance~\cite{xiao2023cpack}. This baseline tests whether a high-performing off-the-shelf retriever already captures narrative structure without task-specific fine-tuning.

\paragraph{Untrained E5-Mistral.}
We include \texttt{intfloat/e5-mistral-7b-instruct} before structural fine-tuning. E5 models are trained with weakly supervised contrastive objectives for general-purpose text embedding~\cite{wang2022text}, and the Mistral-based instruction variant improves text embeddings with large language models~\cite{wang2023improving}. Comparing against this baseline isolates the effect of our structural pseudo-label fine-tuning.

\paragraph{StoryEmb.}
We evaluate \texttt{uhhlt/story-emb}, a narrative-focused embedding model trained for story retrieval and released with the Story Embeddings work~\cite{hatzel2024story}. This is the closest baseline in spirit to our goal, since it is explicitly designed to represent fictional stories according to narrative properties rather than generic topical similarity.

\paragraph{GPT-5.}
We use GPT-5 as a prompt-based judge baseline on tasks that can be evaluated by direct comparison or scoring~\cite{openai2025gpt5}. Unlike embedding baselines, GPT-5 reads the input stories at inference time and produces a judgment directly, so it is a strong but comparatively expensive reference point rather than a scalable retrieval model. We therefore report GPT-5 separately for tasks where a direct LLM decision is available.

\section{Results}

\subsection{Embedding Checkpoint and Baseline Comparison}

Table~\ref{tab:embedding-baseline-comparison} compares the v2 structural MSE checkpoint against the embedding baselines available in the saved evaluation run. We report the pairwise cosine separation gap, retrieval accuracy, ranking nDCG, SemEval 2026 triplet accuracy, synthetic anchor-comparison accuracy, and correlation with SemEval 2022 NAR scores. GPT-5 is included where we have a directly comparable prompt-based result.

\begin{table*}[t]
\centering
\small
\begin{tabular}{lrrrrrrr}
\toprule
Model & Pair gap & Ret. acc. & Rank nDCG & Syn. acc. & S26 acc. & S22 $\rho$ & S22 $r$ \\
\midrule
BGE-base & 0.146 & 0.885 & 0.962 & 0.450 & 0.615 & -0.737 & -0.745 \\
E5-Mistral & 0.162 & 0.970 & 0.975 & 0.450 & 0.575 & -0.692 & -0.624 \\
RoBERTa-base & 0.007 & 0.505 & 0.913 & 0.300 & 0.500 & -0.161 & -0.137 \\
StoryEmb & 0.172 & 0.955 & 0.976 & 0.400 & 0.570 & -0.685 & -0.634 \\
\textbf{E5-Mistral+MSE v2 e1} & \textbf{0.427} & \textbf{0.970} & \textbf{0.978} & 0.450 & \textbf{0.630} & -0.723 & -0.712 \\
GPT-5 / GPT-5.4 & -- & -- & -- & \textbf{0.750} & \textbf{0.720} & -- & -- \\
\bottomrule
\end{tabular}
\caption{Checkpoint and baseline comparison from the evaluation pipeline. Pair gap is the mean cosine similarity of positive pairs minus the mean cosine similarity of negative pairs. Ret. acc. is five-way retrieval accuracy. Syn. acc. is synthetic anchor-comparison accuracy. S26 acc. is SemEval 2026 Task 4 triplet accuracy. S22 $\rho$ and S22 $r$ are Spearman and Pearson correlations between embedding cosine similarity and SemEval 2022 NAR scores. GPT-5 is not an embedding model, so cosine-based and retrieval-indexing metrics are not directly applicable.}
\label{tab:embedding-baseline-comparison}
\end{table*}

The v2 structural MSE checkpoint gives the strongest embedding-based pairwise separation, increasing the positive-negative cosine gap from 0.172 for the strongest baseline to 0.427. It also matches the best retrieval accuracy among embedding models, slightly improves ranking nDCG, and obtains the best embedding-based SemEval 2026 triplet accuracy. The SemEval 2022 news-article correlation remains negative across all embedding models, suggesting that this external NAR score may capture a different or differently oriented notion of narrative similarity than the story-structure signal used for training. GPT-5 gives the strongest available direct-judgment performance on SemEval 2026 and the synthetic anchor task, but those scores require full prompt-time comparison rather than reusable embeddings.

\subsection{LLM Structural Scoring on Movie Remakes}

Table~\ref{tab:movie-llm-scores} reports GPT-5 structural similarity scores on the Movie Remakes evaluation subset. The LLM assigns substantially higher scores to positive pairs than to random negative pairs, indicating that the prompt captures a clear structural distinction between related and unrelated movie narratives.

\begin{table}[t]
\centering
\small
\begin{tabular}{lrrrrr}
\toprule
Pair type & $n$ & Mean & Std. & Min & Max \\
\midrule
Positive & 50 & 7.80 & 2.12 & 2 & 10 \\
Random negative & 49 & 2.27 & 0.76 & 1 & 4 \\
\bottomrule
\end{tabular}
\caption{GPT-5 structural similarity scores on Movie Remakes pairs. Scores are on a 1--10 scale, where larger values indicate greater structural similarity. One negative pair was excluded because it did not receive a valid score.}
\label{tab:movie-llm-scores}
\end{table}

\subsection{Agreement Between GPT-5 and the Structural Scorer}

We further compare GPT-5 scores against our structural similarity model on 1,000 randomly sampled ASQ story pairs. The LLM scores are discrete ratings from 1 to 10, while the structural model produces continuous calibrated scores. We therefore report both Pearson correlation, which measures linear association, and Spearman correlation, which measures rank-level monotonic agreement.

\begin{table}[t]
\centering
\small
\begin{tabular}{lrrr}
\toprule
Comparison & $n$ & Correlation & $p$-value \\
\midrule
Pearson $r$ & 1000 & 0.138 & $1.16\times10^{-5}$ \\
Spearman $\rho$ & 1000 & 0.136 & $1.65\times10^{-5}$ \\
\bottomrule
\end{tabular}
\caption{Correlation between GPT-5 structural similarity ratings and structural model scores on ASQ random story pairs.}
\label{tab:asq-llm-correlation}
\end{table}

The Movie Remakes results suggest that GPT-5 can distinguish positive from random negative story pairs when asked to focus on event progression and narrative structure. However, the ASQ correlation analysis shows only weak agreement between GPT-5 and the calibrated structural scorer on randomly sampled personal narratives. This suggests that LLM judgments and the structural scorer may capture overlapping but non-identical notions of narrative similarity, motivating direct evaluation of trained embeddings rather than assuming LLM scores are a perfect proxy for the target signal.

\section{Conclusion}

We present a framework for learning structure-aware story embeddings using human-calibrated structural similarity scores and pseudo-labeled narrative pairs. The resulting embeddings are designed for scalable retrieval and recommendation settings where direct LLM-based comparison is prohibitively expensive. Future work should expand human evaluation, improve candidate selection for pseudo-labeling, and test the embeddings in downstream narrative RAG and recommendation systems.


\section*{Limitations}
This draft should be updated with final dataset sizes, checkpoint selections, annotation details, and the complete set of quantitative results before submission.

\section*{Ethics Statement}
This work uses narrative data and automatically generated similarity labels. Any public release should document dataset licenses, filtering decisions, model limitations, and possible biases introduced by pseudo-labeling and LLM-based evaluation.

\nocite{placeholder}
\bibliography{references}

\appendix
\section{Appendix}
Additional prompts, hyperparameters, and evaluation tables can be placed here.

\end{document}
